\begin{tcbitemize}[
  raster columns=3, raster equal height=rows,
  blanker, cheat subtitle
]
  \tcbitem
    \begin{tcbitemize}[
      raster columns=1, raster equal height=rows,
      blanker, cheat subtitle
    ]
      \tcbitem
          \tcbsubtitle{Upgrade a Cluster (kubeadm)}
          {\fontsize{6.5}{7.5}\selectfont\sffamily
          Control plane:
          \begin{enumerate}[leftmargin=1.2em, itemsep=0pt, parsep=0pt, topsep=1pt]
            \item Edit \texttt{/etc/apt/sources.list.d/kubernetes.list} — change version to target minor (e.g.\ \texttt{v1.31})
            \item \texttt{apt-mark unhold kubeadm} — allow apt to upgrade it
            \item \texttt{apt update \&\& apt install -y kubeadm=1.X.Y-1.1}
            \item \texttt{kubeadm upgrade plan}
            \item \texttt{kubeadm upgrade apply v1.X.Y}
            \item \texttt{apt install -y kubelet=1.X.Y-1.1 kubectl=1.X.Y-1.1}
            \item \texttt{systemctl daemon-reload \&\& systemctl restart kubelet}
          \end{enumerate}
          Worker nodes:
          \begin{enumerate}[leftmargin=1.2em, itemsep=0pt, parsep=0pt, topsep=1pt]
            \item \texttt{k drain <node> --ignore-daemonsets --delete-emptydir-data}
            \item SSH to node, then: \texttt{kubeadm upgrade node}
            \item \texttt{apt install -y kubelet=1.X.Y-1.1 kubectl=1.X.Y-1.1}
            \item \texttt{systemctl daemon-reload \&\& systemctl restart kubelet}
            \item Exit SSH back to control plane node
            \item \texttt{k uncordon <node>}
          \end{enumerate}}

      \tcbitem
          \tcbsubtitle{Create ServiceAccount + Bind Role}
          {\fontsize{6.5}{7.5}\selectfont\sffamily
          \begin{enumerate}[leftmargin=1.2em, itemsep=0pt, parsep=0pt, topsep=1pt]
            \item \texttt{k create sa deploy-admin -n <ns>}
            \item \texttt{k create clusterrole cr --verb=get --resource=pods}
            \item \texttt{k create clusterrolebinding crb --clusterrole=cr --serviceaccount=<ns>:deploy-admin}
            \item \texttt{k auth can-i get pods --as system:serviceaccount:<ns>:deploy-admin  \# Verify}
          \end{enumerate}}

      \tcbitem
          \tcbsubtitle{Curl the K8s API from Inside a Pod}
          {\fontsize{6.5}{7.5}\selectfont\sffamily
          Every pod gets a ServiceAccount token auto-mounted at \texttt{/var/run/secrets/kubernetes.io/serviceaccount/}.}
          \begin{cheatcodebash}
          # Inside a pod (e.g. k exec -it <pod> -- sh):
          TOKEN=$(cat /var/run/secrets/kubernetes.io/\
          serviceaccount/token)
          curl -k https://kubernetes.default.svc/api/v1/\
          namespaces/default/pods \
            -H "Authorization: Bearer $TOKEN"
          # -k = skip TLS verify (self-signed cluster CA)
          \end{cheatcodebash}

      \tcbitem
          \tcbsubtitle{Create Deployment, Upgrade \& Annotate}
          \begin{cheatcodebash}
          k create deploy nginx --image=nginx:1.24 --replicas=3
          k set image deploy/nginx nginx=nginx:1.25
          k annotate deploy nginx kubernetes.io/change-cause="upgrade to 1.25"
          \end{cheatcodebash}

      \tcbitem
          \tcbsubtitle{Create NetworkPolicy}
          {\fontsize{6.5}{7.5}\selectfont\sffamily
          \begin{enumerate}[leftmargin=1.2em, itemsep=0pt, parsep=0pt, topsep=1pt]
            \item Identify pods via labels: \texttt{k get pods --show-labels}
            \item Write policy with \texttt{podSelector}, \texttt{ingress}/\texttt{egress} rules
            \item \emoji{warning}  An empty \texttt{podSelector: \{\}} selects \textbf{all} pods in the namespace.
            \item Test: \texttt{k exec <pod> -- wget -qO- -T2 <target>}
          \end{enumerate}}
          \begin{cheatcodebash}
          # docs: search "network policies"
          \end{cheatcodebash}

      \tcbitem
          \tcbsubtitle{Custom Columns \& Sorting}
          {\fontsize{6.5}{7.5}\selectfont\sffamily
          List all deployments with name, replicas, and image, sorted by name:}
          \begin{cheatcodebash}
          k get deploy -A -o custom-columns=\
              'NAME:.metadata.name,NS:.metadata.namespace,\
              REPLICAS:.spec.replicas,\
              IMAGE:.spec.template.spec.containers[0].image' \
            --sort-by=.metadata.name
          \end{cheatcodebash}
          \vspace{-0.5ex}
          {\fontsize{6.5}{7.5}\selectfont\sffamily
          Sort pods by restart count:}
          \begin{cheatcodebash}
          k get pods -A -o custom-columns='NAME:.metadata.name,\
             RESTARTS:.status.containerStatuses[0].restartCount' \
             --sort-by='.status.containerStatuses[0].restartCount'
          \end{cheatcodebash}

    \end{tcbitemize}

  \tcbitem
    \begin{tcbitemize}[
      raster columns=1, raster equal height=rows,
      blanker, cheat subtitle
    ]

      \tcbitem
          \tcbsubtitle{etcd Backup \& Restore}
          {\fontsize{6.5}{7.5}\selectfont\sffamily
          \begin{enumerate}[leftmargin=1.2em, itemsep=0pt, parsep=0pt, topsep=1pt]
            \item Check etcd version: \texttt{k -n kube-system exec etcd-* -- etcd --version}
            \item Get certs from \texttt{k describe pod etcd-* -n kube-system}\\
              Map etcd pod flags to \texttt{etcdctl} flags:\\
              \texttt{--advertise-client-urls} → \texttt{--endpoints}\\
              \texttt{--cert-file} → \texttt{--cert}\\
              \texttt{--key-file} → \texttt{--key}\\
              \texttt{--trusted-ca-file} → \texttt{--cacert}
            \item Backup:
          \end{enumerate}}
          \begin{cheatcodebash}
          ETCDCTL_API=3 etcdctl snapshot save /tmp/backup.db \
            --endpoints=https://127.0.0.1:2379 \
            --cacert=/etc/kubernetes/pki/etcd/ca.crt \
            --cert=/etc/kubernetes/pki/etcd/server.crt \
            --key=/etc/kubernetes/pki/etcd/server.key
          \end{cheatcodebash}
          \vspace{-0.5ex}
          {\fontsize{6.5}{7.5}\selectfont\sffamily
          \begin{enumerate}[leftmargin=1.2em, itemsep=0pt, parsep=0pt, topsep=1pt, start=3]
            \item Stop \textbf{all} controlplane components:
          \end{enumerate}}
          \begin{cheatcodebash}
          cd /etc/kubernetes/manifests
          mv *.yaml ..        # kubelet stops all static pods
          watch crictl ps     # wait until ALL containers gone!
          \end{cheatcodebash}
          \vspace{-0.5ex}
          {\fontsize{6.5}{7.5}\selectfont\sffamily
          \begin{enumerate}[leftmargin=1.2em, itemsep=0pt, parsep=0pt, topsep=1pt, start=4]
            \item Restore snapshot (offline, no certs needed):
          \end{enumerate}}
          \begin{cheatcodebash}
          # etcd <3.6: etcdctl snapshot restore
          # etcd >=3.6: etcdutl snapshot restore
          etcdutl snapshot restore /tmp/backup.db \
            --data-dir=/var/lib/etcd-restored
          \end{cheatcodebash}
          \vspace{-0.5ex}
          {\fontsize{6.5}{7.5}\selectfont\sffamily
          \begin{enumerate}[leftmargin=1.2em, itemsep=0pt, parsep=0pt, topsep=1pt, start=5]
            \item Update etcd manifest \texttt{hostPath} to new dir:
          \end{enumerate}}
          \begin{cheatcodebash}
          vi /etc/kubernetes/etcd.yaml
          # Change volumes[name=etcd-data].hostPath.path:
          #   /var/lib/etcd -> /var/lib/etcd-restored
          \end{cheatcodebash}
          \vspace{-0.5ex}
          {\fontsize{6.5}{7.5}\selectfont\sffamily
          \begin{enumerate}[leftmargin=1.2em, itemsep=0pt, parsep=0pt, topsep=1pt, start=6]
            \item Move manifests back, wait for restart:
          \end{enumerate}}
          \begin{cheatcodebash}
          mv /etc/kubernetes/*.yaml /etc/kubernetes/manifests/
          watch crictl ps     # wait for containers to start
          # Can take several minutes!
          \end{cheatcodebash}
          \vspace{-0.5ex}
          {\fontsize{6.5}{7.5}\selectfont\sffamily
          \begin{enumerate}[leftmargin=1.2em, itemsep=0pt, parsep=0pt, topsep=1pt, start=7]
            \item Verify: \texttt{k get pods -A} — cluster state matches snapshot
          \end{enumerate}}
          \begin{cheatcodebash}
          # docs: search "backing up etcd cluster"
          \end{cheatcodebash}

      \tcbitem
          \tcbsubtitle{Create Pod with Secret Mounted}
          \begin{cheatcodebash}
          k create secret generic db-creds \
            --from-literal=user=admin \
            --from-literal=pass=s3cret $do > secret.yaml
          \end{cheatcodebash}
          \vspace{-0.5ex}
          {\fontsize{6.5}{7.5}\selectfont\sffamily
          Then generate pod YAML and add the secret volume:}
          \begin{cheatcodebash}
          k run secret-pod --image=nginx $do > pod.yaml
          # Edit pod.yaml: add volumes + volumeMounts
          \end{cheatcodebash}
          \vspace{-0.5ex}
          \begin{cheatcodeyaml}
          # Add to pod spec:
          volumes:
          - name: secret-vol
            secret:
              secretName: db-creds
          # Add to container:
          volumeMounts:
          - name: secret-vol
            mountPath: /etc/db-creds
            readOnly: true
          \end{cheatcodeyaml}

    \end{tcbitemize}

  \tcbitem
    \begin{tcbitemize}[
      raster columns=1, raster equal height=rows,
      blanker, cheat subtitle
    ]

      \tcbitem
          \tcbsubtitle{Change Service CIDR Range}
          {\fontsize{6.5}{7.5}\selectfont\sffamily
          \begin{enumerate}[leftmargin=1.2em, itemsep=0pt, parsep=0pt, topsep=1pt]
            \item Edit kube-apiserver manifest:
          \end{enumerate}}
          \begin{cheatcodebash}
          vi /etc/kubernetes/manifests/kube-apiserver.yaml
          # Change: --service-cluster-ip-range=10.96.0.0/12
          # To:     --service-cluster-ip-range=10.100.0.0/16
          \end{cheatcodebash}
          \vspace{-0.5ex}
          {\fontsize{6.5}{7.5}\selectfont\sffamily
          \begin{enumerate}[leftmargin=1.2em, itemsep=0pt, parsep=0pt, topsep=1pt, start=2]
            \item Edit kube-controller-manager manifest:
          \end{enumerate}}
          \begin{cheatcodebash}
          vi /etc/kubernetes/manifests/kube-controller-manager.yaml
          # Change: --service-cluster-ip-range=10.100.0.0/16
          # (must match kube-apiserver)
          \end{cheatcodebash}
          \vspace{-0.5ex}
          {\fontsize{6.5}{7.5}\selectfont\sffamily
          \begin{enumerate}[leftmargin=1.2em, itemsep=0pt, parsep=0pt, topsep=1pt, start=3]
            \item Update ServiceCIDR objects (K8s 1.29+):
          \end{enumerate}}
          \begin{cheatcodebash}
          # Create new ServiceCIDR for the new range:
          k apply -f - <<'EOF'
          apiVersion: networking.k8s.io/v1beta1
          kind: ServiceCIDR
          metadata: {name: new-cidr}
          spec:
            cidrs: ["10.100.0.0/16"]
          EOF
          # Delete old ServiceCIDR resource:
          k delete servicecidr <old-cidr-name>
          \end{cheatcodebash}
          \vspace{-0.5ex}
          {\fontsize{6.5}{7.5}\selectfont\sffamily
          \begin{enumerate}[leftmargin=1.2em, itemsep=0pt, parsep=0pt, topsep=1pt, start=4]
            \item Optional: redeploy services to pick up new range:
          \end{enumerate}}
          \begin{cheatcodebash}
          # Delete and recreate services (new ClusterIPs):
          k get svc -A -o yaml > all-svc.yaml
          k delete svc <name> -n <ns>
          k apply -f all-svc.yaml
          \end{cheatcodebash}

      \tcbitem
          \tcbsubtitle{Install Container Runtime}
          {\fontsize{6.5}{7.5}\selectfont\sffamily
          \begin{enumerate}[leftmargin=1.2em, itemsep=0pt, parsep=0pt, topsep=1pt]
            \item \texttt{ssh <worker-node>}
            \item \texttt{sudo -i}
            \item \texttt{dpkg -i /path/to/containerd.io\_<ver>\_amd64.deb}
            \item Configure containerd for systemd cgroup:
          \end{enumerate}}
          \begin{cheatcodebash}
          mkdir -p /etc/containerd
          containerd config default > /etc/containerd/config.toml
          # Set SystemdCgroup = true in config.toml:
          sed -i 's/SystemdCgroup = false/SystemdCgroup = true/' \
            /etc/containerd/config.toml
          \end{cheatcodebash}
          \vspace{-0.5ex}
          {\fontsize{6.5}{7.5}\selectfont\sffamily
          \begin{enumerate}[leftmargin=1.2em, itemsep=0pt, parsep=0pt, topsep=1pt, start=5]
            \item Enable and start the service:
          \end{enumerate}}
          \begin{cheatcodebash}
          systemctl daemon-reload
          systemctl enable --now containerd
          systemctl status containerd   # verify Active: active
          \end{cheatcodebash}
          \vspace{-0.5ex}
          \begin{cheatcodebash}
          # docs: search "container runtimes"
          \end{cheatcodebash}

      \tcbitem
          \tcbsubtitle{Update a Helm Chart}
          {\fontsize{6.5}{7.5}\selectfont\sffamily
          Naming: \textbf{release} = installed instance, \textbf{repo} = chart source, \textbf{chart} = package name.}
          \begin{cheatcodebash}
          # 1. Find release name, namespace, chart
          helm ls -A   # → NAME  NAMESPACE  CHART  ...
          # 2. List configured repos
          helm repo ls
          # 3. Refresh repo index (global, not namespaced)
          helm repo update <repo>
          # 4. List all available versions of a chart
          helm search repo <repo>/<chart> -l
          # 5. Upgrade release to specific version
          helm upgrade <release> <repo>/<chart> -n <ns> --version=x.y.z
          \end{cheatcodebash}
          \vspace{-0.5ex}
          \begin{cheatcodebash}
          # docs: search "helm" (helm.sh/docs)
          \end{cheatcodebash}

    \end{tcbitemize}

\end{tcbitemize}
